\section{Conclusiones}
Se desarrolló e implementó \textbf{UrbanTracker}, un sistema innovador de geolocalización en tiempo real orientado al transporte público urbano, que integra aplicaciones para usuarios, conductores y administradores en una plataforma unificada. El sistema cumple con los objetivos planteados de ofrecer información inmediata y precisa sobre la ubicación de autobuses a los pasajeros, y de facilitar la gestión operativa a los administradores mediante herramientas centralizadas. La arquitectura adoptada -- combinando comunicación en tiempo real (WebSockets/MQTT), backend modular en Spring Boot y visualización con Mapbox -- permitió alcanzar frecuencias de actualización de aproximadamente 5 segundos y mantener latencias generalmente inferiores a 2 segundos en las pruebas realizadas, satisfaciendo así los criterios de desempeño clave. Además, el uso de estándares abiertos y hardware común redujo los costos de implementación, favoreciendo la adopción en contextos con recursos limitados.

En términos de impacto, UrbanTracker contribuye a reducir la incertidumbre de los tiempos de espera para los pasajeros y a mejorar la toma de decisiones operativas para las entidades gestoras del transporte. A futuro, se proyecta extender la plataforma con funcionalidades como estimación de tiempos de llegada, alertas, y análisis históricos de recorridos para apoyo a la planificación. También se plantea profundizar en aspectos de \textbf{seguridad y privacidad}, incluyendo políticas de retención de datos y controles de acceso más finos, y realizar evaluaciones de usabilidad a mayor escala con usuarios reales. Estas mejoras consolidarán a UrbanTracker como una solución robusta, escalable y centrada en el usuario para la gestión de flotas de transporte urbano.
